\subsection{buffer\_pool}\label{buffer_pool}
Das \verb|buffer_pool| Modul verwaltet eine feste Anzahl vorab allozierter
\verb|buffer_t| (Abschnitt \ref{buffer}) und vergibt sie nach dem
Belegt\-/Frei\-Prinzip. Anwender holen sich mit \verb|buffer_pool_get()|
einen freien Puffer und geben ihn nach Gebrauch mit
\verb|buffer_pool_return()| zurück, statt einzelne Puffer wiederholt zu
allozieren und freizugeben. Das vermeidet Heap\-Fragmentierung und macht den
Speicherbedarf zur Laufzeit deterministisch.

\subsubsection*{Zustand}
Der Pool wird durch \verb|buffer_pool_t| (Listing \ref{buffer_pool_t})
beschrieben. \verb|buffer| ist das Feld der verwalteten
\verb|buffer_cnt| Puffer; \verb|type| bezeichnet die Betriebsart des Pools
und \verb|next| den Index des nächsten Elements (negativ, wenn ein linearer
Puffer im Ring verwendet wird). \verb|sbuffer| ist ein zweites, derzeit für
Erweiterungen reserviertes Feld.

\begin{listing}
\begin{minted}{C}
typedef struct buffer_pool_s {
    uint8_t buffer_cnt;  // Anzahl Puffer im Pool
    buffer_t **buffer;   // Feld der Puffer
    buffer_t **sbuffer;  // reserviert
    b_type_e type;       // Betriebsart des Pools
    int8_t next;         // Index des naechsten Elements
} buffer_pool_t;
\end{minted}
\caption{Definition der \texttt{buffer\_pool\_t} Struktur}
\label{buffer_pool_t}
\end{listing}

\subsubsection*{API}
{\sloppy\emergencystretch=3em
\begin{description}
\item[\texttt{buffer\_pool\_new(lcnt, charCnt, line\_type, type)}] Erzeugt
  einen Pool aus \verb|lcnt| Puffern zu je \verb|charCnt| Bytes. \verb|line_type|
  legt den Typ der einzelnen Puffer fest (\verb|LINEAR| oder \verb|RING|),
  \verb|type| die Betriebsart des Pools. Schlägt eine Teil\-Allokation fehl,
  werden bereits angelegte Puffer wieder freigegeben und \verb|NULL|
  zurückgeliefert.
\item[\texttt{buffer\_pool\_free(pool)}] Gibt alle Puffer und den Pool selbst
  frei.
\item[\texttt{buffer\_pool\_get(bp)}] Liefert den ersten freien Puffer, markiert
  ihn als \verb|BUFFER_USED| und gibt ihn zurück; \verb|NULL|, wenn kein Puffer
  frei ist.
\item[\texttt{buffer\_pool\_return(bp, buffer)}] Gibt einen Puffer an den Pool
  zurück. Geprüft werden der Index\-Bereich, die Zugehörigkeit zum Pool
  (Pointer\-Vergleich) und dass der Puffer nicht mehr belegt ist; danach wird er
  mit \verb|buffer_reset()| zurückgesetzt. Ein noch als \verb|BUFFER_USED|
  markierter Puffer wird mit \verb|EM_ERR| abgelehnt -- der Inhalt muss also
  zuvor verbraucht (z.\,B.\ per \verb|buffer_get()|) worden sein.
\item[\texttt{buffer\_pool\_print(bp, title)}] Gibt die Puffer des Pools zu
  Debug\-Zwecken aus.
\end{description}
\par}
